\cotefile{157-5}

\scanpage{573}

\mspage{482}

ce qui signifie que $x_\alpha$ n'appartient à aucun des $\deleted{\unsure{S_{\alpha}}}\underline S_\beta={}_{x_\beta}M$, pour $\beta\in J$. Donc $\underline S$ n'est pas réunion d'une petite sous-famille de la famille des $ {}_{x_\alpha}M$, donc en vertu de $\unsure{4.16.7.1|4.6.7.1}$, $\underline S$ n'est pas accessible dans $M$.

Inversement, soit $\underline S$ une telle sous-catégorie de $M$. Considérons un bon ordre sur $\mathrm{Ob}\,\underline S$ \inserted{(objets $s_\alpha$)}, et construisons la suite des $x_\alpha$ \inserted{$(\geq s_\alpha)$} par récurrence transfinie.

% slanted
\margin{\uline{NB} On a $\operatorname{card}\underline S\notin U$, \deleted{\unsure{sin}}sinon $\underline S$ serait petite, donc accessible dans $M$, contrairement à l'hypothèse.}

1) On choisit $\deleted{x_0\in\mathrm{Ob}\,\underline S}$ $x_0={}$ premier élément de $\mathrm{Ob}\,S$ (existe, car $\mathrm{Ob}\,\underline S\ne\varnothing$ i.e. $\underline S\ne\varnothing$, ou sinon $\underline S$ serait accessible dans $M$).

2) Si les $x_\beta$ ($\beta<\alpha$) sont construits, on prend pour $x_\alpha$ le premier élément de $\underline S$ \inserted{\unsure{majorant} les $x_\beta$ et} qui n'appartient pas à la sous-catégorie fermée de $M$ engendrée par les $x_\beta$ ($\beta<\alpha$). [Cette sous-catégorie fermée est $\subset\underline S$, et $\ne\underline S$ car $\underline S$ n'est pas accessible dans $M$] On a donc bien

\[
\operatorname{Hom}(x_\beta,x_\alpha)=\varnothing\quad\text{si }\beta<\alpha
\]

% archive sheet number: 572
\scanpage{574}

\mspage{483}

Ainsi, la récurrence transfinie se poursuit sur $I$ \deleted{\unsure{tout}} entier. \deleted{Par construction} On voit par récurrence transfinie qu'on a, pour \unsure{tt} $\alpha\in I$

\[s_\alpha\leq x_\alpha\]

(car $x_\alpha>x_\beta\geq s_\beta$ pour \deleted{\unsure{tout}} $\beta<\alpha$, donc $x_\alpha>s_\beta$ pour \unsure{tt} $\beta<\alpha$, \deleted{donc $x_\alpha\geq s_\alpha$} ce qui équivaut à $x_\alpha\geq s_\alpha$). De plus, on trouve \deleted{\unsure{qu'}}\unsure{aussi}, par récurrence \deleted{\unsure{ind}}transfinie, que

\[
\begin{aligned}
\beta<\alpha\ \Longrightarrow\ s_\beta\ &\text{appartient à la sous-catégorie}\\
 &\text{fermée }\inserted{\underline S_\alpha}\text{ de }M\text{ engendrée}\\
 &\text{par les }x_{\beta'}\text{ avec }\beta'<\alpha .
\end{aligned}\tag*{\makebox[0pt][r]{$(*)$\hspace{.90\linewidth}}}
\]

puisque par construction $s_\alpha$ est le \uline{premier} des éléments de $\underline S$ \deleted{tel que} qui n'est \uline{pas} dans \unsure{la dite|ladite} sous-catégorie \deleted{\unsure{de}} $\underline S_\alpha$ de $M$.

Cela montre que \deleted{\unsure{tout}} $s_\beta$ appartient à la sous-catégorie \inserted{fermée} $ {}_{\underline x}M$ de $M$ engendrée par \deleted{\unsure{tous}} \inserted{la famille} de \uline{tous} les \deleted{$\unsure{x}$} $x_\alpha$, \deleted{donc} on trouve que $\underline S\subset{}_{\underline x}M\subset\underline S$ donc $\underline S={}_{\underline x}M$, qed.

\uline{Proposition} 4.16.7.5. Pour un univers $\unsure{\mathcal U|U}$ donné (contenant $\unsure{\mathbb N|N}$ \inserted{comme élément}), les conditions suivantes sont équivalentes:

1) Il existe une catégorie $\mathcal U$-accessible $M$, \deleted{\unsure{et}}\inserted{\unsure{ayant}} une \deleted{sous-}catégorie fermée non accessible

% archive sheet number: 573
\scanpage{575}

\mspage{484}

\unsure{dans} $M$.

b) Il existe une catégorie $C\in\deleted{\unsure{\mathrm{Cat}}}\unsure{\mathcal U|U}$, et une suite transfinie d'éléments de

\[
C^\wedge\overset{\mathrm{déf}}{=}\underline{\operatorname{Hom}}(C,\deleted{\unsure{\mathcal U}}\mathrm{Ens}_{\mathcal U}),
\]

$\underline x=(x_\alpha)_{\alpha\in I}$, indexée par l'ens de\deleted{\unsure{s}} tous les ordinaux $\alpha\in\mathcal U$, et telle que

\[\operatorname{Hom}(x_\alpha,x_\beta)=\varnothing\qquad\text{si }\alpha<\beta .\]

\uline{Dém} On a $\mathrm{b)}\Rightarrow\mathrm{a)}$ en prenant $M=C^\wedge$, en vertu de la proposition précédente. On voit de \unsure{suite|même|mê} $\mathrm{a)}\Rightarrow\mathrm{b)}$, en plongeant $M$ dans une catégorie $C^\wedge$ ($C$ étant une petite sous-catégorie \unsure{str.} génératrice de $M$).

Je n'ai pas été capable$^{*}$%
% slanted
\margin{Mais si finalement, cf 4.16.7.10 plus bas} de construire une suite $(x_\alpha)_{\alpha\in I}$ comme exigée dans cette proposition, ni de prouver qu'une telle suite n'existe pas, même en un cas particulier tel que $M=k\text{-}\unsure{\mathrm{alg\,comm}|\mathrm{dg\,comm}|\mathrm{dgcom}}$, $k$ un anneau commutatif. \deleted{\unsure{Même en|Notons que}} Je n'ai pas

% archive sheet number: 574
\scanpage{576}

\mspage{485}

$=$ prouvé qu'on ne puisse trouver une suite transfinie de \uline{corps} $k_\alpha$, ayant les propriétés dites.

A fortiori, je n'ai pu tirer au clair la question si une sous-catégorie $\underline L$ stable d'une catégorie $M$ accessible est tjrs%
% slanted
\margin{C'est faux cf plus bas} accessible dans $M$. Il se pourrait qu'il n'y ait pas de tel cas. Il se pourrait aussi qu'il y en ait, même avec une sous-catégorie fermée (laquelle est stable par \unsure{tt} type de \uline{lim} représentable dans $M$, sauf la \uline{lim} sur une catégorie d'indices \unsure{vides|vide}\ldots).

Pour les sous-catégories \uline{ouvertes} je rappelle que j'ai construit (\unsure{dès|dans} l'appendice) un exemple d'une \unsure{telles|telle} \unsure{sous-catégories|sous-catégorie} non accessible dans $M$. Au sujet des sous-catégories ouvertes, je signale ici l'énoncé (quasi tautologique) :

% archive sheet number: 575
\scanpage{577}

\mspage{486}

% slanted
\margin{Cette proposition et son \unsure{corollaire} sont \unsure{contenus|contenues} dans 4.16.7.8 ci-dessous, il n'y a pas lieu d'en faire un énoncé \unsure{séparé|à part}.}

\uline{Proposition} 4.16.7.6. Soient $M$ une catégorie accessible, $\underline S$ une sous-catégorie \uline{ouverte} de $M$. Pour que $\underline S$ soit accessible dans $M$, il \unsure{f} et s. qu'elle soit $\underline L$-stable dans $M$, i.e. qu'il existe un cardinal $\pi_0$ tel que $M$ satisfasse à $L_{\pi_0}$, et que $\underline S$ soit stable dans $M$ par \uline{lim} de type $L_{\pi_0}$.

De façon plus précise, si cette condition est satisfaite, alors pour tt cardinal $\pi$ tel que $\pi\geq\pi_0$ et que $\pi$ soit utile pour $M$, $\underline S$ est $\pi$-accessible dans $M$. \inserted{Ou encore: si $\pi$ est un cardinal utile pour $M$, alors $\underline S$ est $\pi$-acc. dans $M$ sss il est $L_\pi$-\unsure{stable|st-stable} dans $M$.} Cette dernière assertion résulte du \inserted{\unsure{aussitôt du}} corollaire suivant, lui-même tautologique:

\uline{Corollaire} 4.16.7.7. Soit $\pi$ un cardinal comme \unsure{désigné|requis} ci-dessus, \deleted{et soit} soit $\underline S_{(\pi)}=\underline S\cap M_\pi$, et soit $x\in\mathrm{Ob}\,M$ et

\[
x\simeq\varinjlim_{i\in I}x_i
\quad\left|\begin{aligned}
 &x_i\in\mathrm{Ob}\,M_\pi\\
 &I\ \text{\unsure{grd|gd} devant }\pi
\end{aligned}\right.
\]

une représentation de $x$ comme \deleted{$\varinjlim$} image essentielle d'un objet de $\unsure{\mathrm{Ind}|\mathit{Ind}}_\pi M_\pi$.

\[\text{Alors }\ x\in\mathrm{Ob}\,\underline S\ \Longleftrightarrow\ x_i\in\mathrm{Ob}\,\underline S_{(\pi)}\qquad\forall i\in I .\]

% archive sheet number: 576
\scanpage{578}

\mspage{487}

J'ai finalement tiré au clair la question de l'accessibilité des sous-catégories \deleted{\unsure{fermées}} \deleted{\unsure{et|de}} des sous-catégories fermées d'une catégorie accessible :

\uline{Proposition} $4.16.\unsure{7}.8$. \unsure{Soient} $M$ une \unsure{catégories|catégorie} accessible, $\underline S$ une sous-catégorie \deleted{fermée} de $M$, $\underline T$ la sous-catégorie ouverte complémentaire de $\underline S$, $\pi$ un cardinal infini utile pour $M$. Les conditions suivantes sont équivalentes :

a) $\underline S$ est $\pi$-accessible dans $M$

a$'$) $\underline S$ est le fermé engendré par $\underline S_{(\pi)}\overset{\mathrm{déf}}{=}\underline S\cap M_\pi$, i.e. pour $x\in\mathrm{Ob}\,M$, on a $x\in\mathrm{Ob}\,\underline S$ sss $\exists\,\xi\in\mathrm{Ob}\,\underline S_{(\pi)}$ tel que $\operatorname{Hom}_M(\xi,x)\ne\varnothing$.

b) $\underline T$ est $\pi$-accessible dans $M$.

b$'$) $\underline T$ est $\underline L_\pi$-stable dans $M$.

Notons \deleted{\unsure{tout}} de suite le \deleted{\unsure{cor}}

\uline{Corollaire} 4.16.7.9. Pour $M,\underline S,\underline T$ comme dessus, conditions équivalentes :

a) $\underline S$ est accessible dans $M$.

a$'$) $\exists$ une \uline{petite} partie $A$ de $\mathrm{Ob}\,M$, telle que $\underline S$ soit la sous-catégorie fermée $ {}_A M$ engendrée par $A$, i.e. $\underline S=\bigcup_{\xi\in A}{}_\xi M$, où $ {}_a M$ est \deleted{l'image \unsure{dans|d'un}} la sous-catégorie \unsure{fermée|pleine} de $M$ formée des $x\in\mathrm{Ob}\,M$ tels que $\operatorname{Hom}(\xi,x)\ne\varnothing$.

% archive sheet number: 577
\scanpage{579}

\mspage{488}

b) $T$ est accessible dans $M$

b$'$) $T$ est $\underline L$-stable dans $M$.

\uline{Corollaire} 4.16.7.10. Soit $M$ une catégorie accessible. Les conditions suivantes sont équivalentes

a) Il existe une sous-catégorie fermée $\deleted{\unsure{\underline S|S}}$ de $M$, non accessible dans $M$.

b) Il existe une sous-catégorie ouverte $\deleted{\unsure{\underline T|\underline S}}$ de $M$, non accessible dans $M$.

% slanted
\margin{Exemple où ces conditions sont satisfaites: $M=k\text{-algèbres}$, $k$ un corps \unsure{comm.} infini, $\unsure{\mathcal U|U}$ étant le plus petit univers tel que $\unsure{\mathbb N|N}\in\mathcal U$.}

c) Il existe une suite transfinie \unsure{$x$}

\[\underline x=(x_\alpha)_{\alpha\in I}\]

d'objets de $M$, indexée par l'ens $I$ des ordinaux de l'Univers de référence $\mathcal U$, telle que l'on ait

\[\operatorname{Hom}(x_\alpha,x_\beta)=\varnothing\qquad\text{pour }\alpha<\beta .\]

En effet, l'équivalence de a) et b) est claire par le corollaire précédent, \deleted{\unsure{la précédente|ce qui précède}} et celle de a) et c) n'est autre que 4.16.7.4.

\uline{Démonstration de} 4.16.7.8. Nous aurons à utiliser le

\uline{Lemme} 4.16.7.11. Soient $M$ une catégorie accessible, $\pi$ un cardinal utile pour $M$, $\underline S$ une sous-catégorie pleine de $M$, $\underline S_{(\pi)}\overset{\mathrm{déf}}{=}\underline S\cap M_\pi$, d'où un foncteur pleinement fidèle

\[\varphi:\mathrm{Ind}_\pi(\underline S_{(\pi)})\longrightarrow M .\]

% archive sheet number: 578
\scanpage{580}

\mspage{489}

a) Soit $x\in\mathrm{Ob}\,M$. Pour que $x$ appartienne à l'image essentielle du foncteur $\unsure{\varphi|\psi}$, il f. et s. que la sous-catégorie pleine $\underline S_{(\pi)/x}$ de $\deleted{\unsure{M_\pi}}\deleted{\unsure{M_\pi}}$ la catégorie $M_{\pi/x}$ (qui est filtrante et $=$ gde devant $\pi$) soit cofinale (donc aussi gde devant $\pi$).

b) Pour que $\underline S$ soit $\pi$-accessible dans $M$, il faut et il suffit que pour tt $x\in\mathrm{Ob}\,M$, on ait

\[x\in\mathrm{Ob}\,\underline S\quad\Longleftrightarrow\quad\underline S_{(\pi)/x}\ \text{est cofinale dans }M_{\pi/x},\]

\inserted{[ou encore, que $\underline S$ satisfasse à $L_\pi$, et que pour tt $x$ de $M$, on ait $x\in\mathrm{Ob}\,\underline S\Longrightarrow\underline S_{(\pi)/x}$ cofinal dans $M_{\pi/x}$]}

\uline{Démonstration}. La partie b) est une conséquence tautologique de a), que je vais prouver.

\deleted{\unsure{Si}}\inserted{Dire que} $x$ \deleted{\unsure{est}} est dans l'image essentielle de $\underline S$, signifie que \inserted{$(\unsure{\alpha|*})\ x\simeq\varinjlim_{i\in I}x_i$}, avec $I$ \unsure{grd|gd} devant $\pi$ et les $x_i$ dans $\underline S_{(\pi)}$. On sait que \deleted{\unsure{de}} \unsure{ceci|ces|ça} \unsure{implique|impliquent} ($I$ étant filtrant, et $M$ étant équivalent à $\mathrm{Ind}_\pi(M)$ puisque $\pi$ utile pour $M$) que $I\longrightarrow M_{\pi/x}$ est cofinale (4.4.2 page 79). Or ce foncteur se factorise par le foncteur d'inclusion pleinement fidèle $\underline S_{(\pi)/x}\unsure{\lhook\joinrel\relbar|\lhook\joinrel\longrightarrow} M_{\pi/x}$, lequel est cofinal \unsure{sss|ssi} il satisfait la condition $\mathrm{F}1$. Il est immédiat que cette condition est satisfaite, quand elle est satisfaite pour le foncteur composé en question, ce qui est bien le cas

% archive sheet number: 579
\scanpage{581}

\mspage{490}

par hypothèse.

Inversement, supposons que $\underline S_{(\pi)/x}$ soit cofinal dans $M_{\pi/x}$. Comme $\deleted{\unsure{M}}$ on a

\[x\overset{\sim}{\leftarrow}\underset{M_{\pi/x}}{\unsure{\underline{\lim}|\varinjlim}}\,x_i,\]

$\pi$ étant utile pour $M$, on aura \inserted{aussi} par cofinalité

\[x\overset{\sim}{\leftarrow}\underset{\underline S_{(\pi)/x}}{\varinjlim}\,x_i,\]

or $M_{\pi/x}$ étant \unsure{grde|gde} devant $\pi$, il en est de $=$ de la sous-catégorie \deleted{\unsure{filtr.}} pleine cofinale $\underline S_{(\pi)/x}$. Donc $x$ est dans l'image essentielle de $\mathrm{Ind}_\pi\underline S_{(\pi)}$, q.e.d

\uline{Lemme} 4.16.7.12. Soient $I$ une catégorie \inserted{filtrante}, $J$ une sous-catégorie fermée, $K$ la sous-catégorie ouverte complémentaire. On a alors les équivalences

\diagram{\begin{tikzcd}[ampersand replacement=\&,column sep=large]
J\text{ cofinal}\arrow[r,Leftrightarrow]\arrow[d,Leftrightarrow] \& K\text{ non cofinal}\arrow[d,Leftrightarrow] \\
J\ne\varnothing \& K\ne I .
\end{tikzcd}}

Démonstration immédiate.

Je reviens à la démonstration de 4.16.7.8. On applique \unsure{les|le} \deleted{deux} \unsure{lemmes|lemme} 4.16.7.11 \deleted{\unsure{aux deux|aux de}} à $M,\pi$ et aux deux

% archive sheet number: 580
\scanpage{582}

\mspage{491}

sous-catégories $\underline S,\underline T$ de $M$, en notant que pour $x\in\mathrm{Ob}\,M$, $\underline S_{(\pi)/x}$ est une \inserted{ss-cat.} fermée de $\underline S_{(\pi)}$, dont la complémentaire ouverte est $\underline T_{(\pi)/x}$, de sorte qu'on peut appliquer le lemme 4.16.7.12. On trouve

\diagram{\begin{tikzcd}[ampersand replacement=\&,column sep=2pt,row sep=6pt,font=\scriptsize,scale=.78,transform shape]
\underline S\ \pi\text{-accessible}\&\overset{4.16.7.11}{\Longleftrightarrow}\&\forall x\in\mathrm{Ob}\,M,\ \inserted{\text{on a}}\&[x\in\mathrm{Ob}\,\underline S\&\Longleftrightarrow\&\underline S_{(\pi)/x}\text{ cof. dans }M_{\pi/x}]\\
\&\Longleftrightarrow\&\forall x\in\mathrm{Ob}\,M,\text{ on a }\&\deleted{\text{\unsure{et on}}}[\underbrace{x\notin\mathrm{Ob}\,\underline S}\&\Longleftrightarrow\&\underbrace{\underline S_{(\pi)/x}\text{ non cof. }\inserted{\text{dans }M_{\pi/x}}}]\\
\&\&\text{par }4.16.7.12\arrow[r]\arrow[rrr,bend right=10]\&(\text{i.e. }x\in\mathrm{Ob}\,\underline T)\&\&\substack{(\text{i.e. }\underline T_{(\pi)/x}\\\text{cof. dans }M_{\pi/x})}\\
\&\overset{4.16.7.11}{\Longleftrightarrow}\&\underline T\text{ est }\pi\text{-accessible.}\&\&\&
\end{tikzcd}}

On a donc l'équivalence de a) et b). D'autre part, $\underline S_{(\pi)/x}$ cofinal dans $M_{\pi/x}$ signifie \unsure{simplem\textsuperscript{t}|simplement} \unsure{que} $\underline S_{(\pi)/x}\ne\varnothing$, \deleted{donc} (4.16.7.12), donc \deleted{\unsure{on}} $\underline S$ $\pi$-accessible signifie que $\underline S$ est formée des $x\in\mathrm{Ob}\,M$ tels que $\underline S_{(\pi)/x}\ne\varnothing$, ce qui signifie aussi que $\underline S$ est \inserted{le fermé dans $M$} engendré par $\underline S_{(\pi)}$, d'où $a\Longleftrightarrow a'$. D'autre part,

\deleted{$\underline T_{(\pi)/x}$ cofinal dans $M_{\pi/x}$ signifie que $\underline T_{(\pi)/x}=M_{\pi/x}$, donc $\underline T$ est $\pi$-accessible dans $M$ sss elle est formée des $x\in\mathrm{Ob}\,M$ tels que $\underline T_{(\pi)/x}={}$}

L'équivalence $b\Longleftrightarrow b'$ n'est autre que $4.6.7.6$\unsure{ --|.} (dernière formulation).

% archive sheet number: 581
\scanpage{583}

\mspage{492}

4.17. Retour sur les \unsure{paratopos}, et sur les opérations $M\boxtimes N$, $\underline{\operatorname{Hom}}_!(M,N)$.

% archive sheet number: 582
\scanpage{584}

$\overline{\underline{\mathrm{XVIII}}}$

\unsure{Catégories|Catégorie} accessibles

\deleted{Modèles} \deleted{(4)}

% archive sheet number: 583
\scanpage{585}

\handwritten{$H_{\mathrm D}^{\unsure{\bullet|\cdot}}(f,\unsure{\xi|\zeta|g})$}

% letterhead
\noindent\begin{minipage}[t]{.48\linewidth}\centering
MINISTERE DE L'EDUCATION NATIONALE\\
ACADEMIE D'AIX-MARSEILLE\\
\rule{1.2cm}{0.3pt}\\
Inspection Académique\\
de Vaucluse\\
\rule{1.2cm}{0.3pt}\\
Tél. : 90 82 99 76\\
Référence : Secrétariat n\textsuperscript{o}
\end{minipage}\hfill
% letterhead
\begin{minipage}[t]{.48\linewidth}
AVIGNON, le\\[1ex]
L'Inspecteur d'académie,\\
directeur des services départementaux\\
de l'Education Nationale de Vaucluse,\\[1ex]
à M
\end{minipage}

\margin{\handwritten{\tikz[baseline=(p.base)]{\node[draw,rounded corners=7pt,inner sep=2pt] (p) {\deleted{\uuline{Perplexités} !}};}}}

% margin note written sideways
\margin{\handwritten{dérivateurs (yoga, constructions) $\big\{$}}

\handwritten{\textcircled{1} \deleted{$W$ \unsure{entraîne|arbitraire} tout Dérivateurs, et le pb d'existence de $u_!,u_*$ ? \unsure{Bidérivateurs|Prédérivateurs}, dér. \unsure{à|;} \unsure{à droite|droite} \ldots}}

\margin{\handwritten{\unsure{Dire que} se prolonge aux topos}}

\handwritten{\inserted{Développer une \uline{\unsure{algèbre} des dérivateurs}}}

\handwritten{\textcircled{2} Critique d'\unsure{Anderson|Andersen}. \unsure{Rattrapages}}

\handwritten{\textcircled{\deleted{\unsure{2}}3} Théorèmes de factorisation \deleted{\unsure{Classes}} Parties spéciales, \uline{\unsure{conjectures|cogénérateurs}}}

\handwritten{\textcircled{4} Théorèmes de stabilité pour structures de Quillen.}

\handwritten{\noindent\makebox[0pt][l]{\raisebox{.45ex}{\rule{.98\linewidth}{.35pt}}}\deleted{\textcircled{5} Catégories et Topos}}

% margin note written sideways
\margin{\handwritten{Cat comme une cat. de modèles $\big\{$}}

\handwritten{\textcircled{\deleted{\unsure{5}}6} Localiseurs dans Cat. Dérivateur associé à un localiseur de Cat, et localiseur associé à un dérivateur.}

\handwritten{\textcircled{\deleted{\unsure{6}}7} L'opération $u_!$ pour $\mathrm{Hot}_W$}

\handwritten{\textcircled{\deleted{\unsure{7}}8} Foncteurs $\underline{\unsure{\mathrm{Ch}|\mathrm{Ob}}}_\infty$, $\underline{\mathrm{Ch}}$, $\underline{\mathrm{Ch}}_\infty^W$, $\underline{\mathrm{Ch}}^W$ et \unsure{astrologie} $1/2$-simpliciale.}

\handwritten{\noindent\rule{.55\linewidth}{0.4pt}}

% margin note written sideways
\margin{\handwritten{formalisme calcul dans les dérivateurs\ldots\unsure{.} $\big\{$}}

\handwritten{\deleted{\unsure{8}}\textcircled{9} L'opération $\underline{\operatorname{Hom}}$ externe dans un dérivateur.}

% archive sheet number: 584
\scanpage{586}

% One oblique stroke crosses the entire draft.
\deleted{\[
F(x',y')\overset{\text{épi}}{\unsure{\longrightarrow|-}}F(x,y)\mathbin{\times}_{G(x,y)}G(x',y')
\qquad\begin{aligned}x&\unsure{\longrightarrow|-}x'\\y&\unsure{\longrightarrow|-}y'\end{aligned}
\]}

\deleted{\[
\begin{aligned}
(f:F\unsure{\longrightarrow|-}G)&\in(\unsure{\Phi|\underline{\varphi}|\mathcal L}\mathbin{\unsure{\otimes|\cup}}\Phi')_*\\
\text{sss}\quad \forall\quad u:x&\longrightarrow x'\qquad\in\Phi\\
v:y&\longrightarrow y'\qquad\in\Phi'
\end{aligned}
\]}

\deleted{\[\text{l'application}\quad
\boxed{F(x',y')\longrightarrow F(x,y)\mathbin{\times}_{G(x,y)}G(x',y')}\quad\text{surjective}
\]}

\deleted{Ce serait \unsure{le|la} TF, et il faudrait que ce soit contenu dans $W\mathbin{\widehat{\boxtimes}}W'$ -- est-ce raisonnable ?}

\deleted{\[
\left\{\begin{aligned}
M\boxtimes \unsure{M'|N'|\mathcal M'} &\qquad \unsure{\mathbf W|W}\mathbin{\widehat{\boxtimes}}\mathbf W'\\
\underline{\operatorname{Hom}}_{!\,\mathrm{loc}}(M,M') &\qquad\text{avec }\deleted{\unsure{\mathit H}}(W,W')
\end{aligned}\right.
\]}

\deleted{\[
M\boxtimes M'\boxtimes I^\wedge\simeq\underline{\operatorname{Hom}}(I^{\unsure{\mathrm o|\circ}},M\boxtimes M')
\]}

\deleted{\[
\begin{aligned}
\underline{\operatorname{Hom}}_{!\,\mathrm{loc}}(M,M')\boxtimes I^\wedge
 &\simeq\underline{\operatorname{Hom}}(I^{\mathrm o},\underline{\operatorname{Hom}}_{!\,\mathrm{loc}}(M,M'))\\
 &\simeq\underline{\operatorname{Hom}}_{!\,\mathrm{loc}}(I^{\mathrm o\wedge},\operatorname{Hom}_{!\,\mathrm{loc}}(M,M'))\\
 &\simeq\operatorname{Hom}_{!\,\mathrm{loc}}(\underset{\substack{\rule{.35pt}{.46ex}\rule{3.2em}{.35pt}\rule{.35pt}{.46ex}\\[-.2ex]M(I^{\mathrm o})\unsure{,}}}{I^{\mathrm o\wedge}\boxtimes M},M')
\end{aligned}
\]}

% archive sheet number: 585
\scanpage{587}

\textcircled{1} Cat. accessibles, cardinaux utiles, \inserted{\unsure{filtr. cat.|filtr. card.|fibr. cat.}} foncteurs accessibles. /

\textcircled{2} Cas des petites catégories.

\deleted{2 -- Parties \unsure{acc.}}

Foncteurs non accessibles $\mathrm{Ens}\to\mathrm{Ens}$

\textcircled{3} Parties accessibles; \deleted{\unsure{th.}} \unsure{définition|définitions}, \unsure{structures|structure}, th. de stabilité. \deleted{\unsure{A}} Cas des Parties $L_{\unsure{1|\gamma|g}}$-acc. \inserted{(en particulier $\unsure{x|\alpha|L}$ $L_\pi$-acc.)}

\textcircled{4} \deleted{\unsure{Parties}} Cas des ss-catég. ouvertes ou fermées. Question pas clarifiée (\unsure{question} d'univers)

\textcircled{5.} Stabilité des catégories accessibles.

\deleted{\textcircled{\unsure{6}}}\textcircled{6} Questions ouvertes.

\textcircled{7} Cas des paratopos, opérations $M\boxtimes N$, $\underline{\operatorname{Hom}}_!(M,N)$, et \unsure{évt.} $\unsure{C|\mathcal C}\mathbin{\overset{\pi}{\boxtimes}}D$,

\[\underline{\operatorname{Hom}}_{!\pi}(C,N)\approxeq\underline{\operatorname{Hom}}_!(\underline{\mathrm{Ind}}_\pi(C),N).\]

-- \textcircled{8} \deleted{\textcircled{\unsure{8}} Cas} Cas où $\deleted{\unsure{M}}N=I^\wedge$, où $N$ un topos.

\noindent\rule{.27\linewidth}{0.4pt}

\textcircled{1} \uline{Th. de factorisation}.

$\cdot$ Parties $Q$-accessibles (\unsure{ou|dites} \unsure{spéciale|spéciales}) de une \inserted{\deleted{paratopos}} catégorie ``spéciale'' (p. ex. un paratopos)

\textcircled{2} \deleted{Questions} Conjectures sur \unsure{caractérisation|cogénération} des dites parties --

\textcircled{3} Théorèmes de stabilité pour structures de Quillen spéciales.

% archive sheet number: 586
